\section{What the editor is}
\label{sec:overview}

LaTeXEditor is a desktop application. It does not upload your documents
anywhere, and it does not typeset in the cloud: it drives the \LaTeX{}
distribution already installed on your machine, and everything it needs to
render --- the code editor, the PDF viewer --- is bundled into the app rather
than fetched from a CDN. Open it on a plane and it behaves exactly as it does
at your desk.

\subsection{The window}

The layout is four regions, which this diagram shows in the positions they
actually occupy:

\begin{center}
\setlength{\fboxsep}{6pt}
\begin{tabular}{c}
  \fbox{\parbox{0.82\textwidth}{\centering \textbf{Toolbar} \\ \small
        document name \quad$\cdot$\quad compiler \quad$\cdot$\quad Typeset}} \\[4pt]
  \begin{tabular}{@{}c@{\hspace{4pt}}c@{\hspace{4pt}}c@{}}
    \fbox{\parbox[t][2.4cm]{0.22\textwidth}{\centering \textbf{Sidebar} \\[3pt] \small
          Files \\ Components \\ AI \\ Git \\ Outline}} &
    \fbox{\parbox[t][2.4cm]{0.30\textwidth}{\centering \textbf{Editor} \\[3pt] \small
          your \texttt{.tex} source}} &
    \fbox{\parbox[t][2.4cm]{0.24\textwidth}{\centering \textbf{Preview} \\[3pt] \small
          the typeset PDF}} \\
  \end{tabular} \\[4pt]
  \fbox{\parbox{0.82\textwidth}{\centering \textbf{Issues} \\ \small
        errors and warnings from the last run}} \\
\end{tabular}
\end{center}

Every divider between those regions can be dragged, and the sidebar, preview
and issues panel can each be hidden --- from the command palette, or by
clicking the \ui{$\times$} on the issues header. The editor is given a minimum
width that the preview is not allowed to squeeze past, so no amount of dragging
or window-resizing can collapse the thing you are writing in.

\subsection{What makes it different}

Three choices are worth knowing about up front, because they shape how the rest
of the tool behaves.

\begin{description}[style=nextline, leftmargin=0pt, labelindent=0pt, labelwidth=0pt, font=\normalfont\itshape]

  \item[The AI edits the document.]
  Most editors bolt on a chat window that produces text you then copy by hand.
  Here you press \kbd{Cmd+K} in the paragraph you want changed, say what you
  want, and get a diff that tints the words that actually differ. Accept and it
  is written back into exactly the range it came from.
  Section~\ref{sec:ai} is devoted to this.

  \item[It knows your project, not just your file.]
  When the assistant generates a table, it is told which packages your preamble
  loads; when it writes a citation, it is given the keys that actually exist in
  your \texttt{.bib}; when it fixes an error, it is handed the error. You are
  never asked to paste a log.

  \item[It typesets locally, once.]
  A single pass of \texttt{pdflatex} (or \texttt{xelatex}, or
  \texttt{lualatex}) per press. This is fast and predictable, and it has one
  consequence you must know about --- see Section~\ref{sec:passes}.

\end{description}

\subsection{Requirements}

You need a \LaTeX{} distribution: \textbf{TeX Live} on Linux and macOS,
\textbf{MacTeX} on macOS, or \textbf{MiKTeX} on Windows. On launch the editor
looks for \texttt{pdflatex}, \texttt{xelatex}, \texttt{lualatex} and
\texttt{latexmk}, and lists whichever it finds in the toolbar's compiler menu.

It does not rely on your \texttt{PATH} alone, which matters more than it
sounds: an app launched from the Dock or Finder inherits a minimal environment
--- roughly \path{/usr/bin:/bin:/usr/sbin:/sbin} --- so a TeX installation that
works perfectly in a terminal can be invisible to a GUI app. The editor
therefore also checks the standard install locations directly
(\path{/Library/TeX/texbin}, Homebrew, MacPorts, per-user \path{~/.local/bin},
and versioned TeX Live trees) and passes the compiler's own directory down to
it, since \texttt{latexmk} shells out to \texttt{pdflatex} in turn.

If no compiler is found, the editor says so in a banner with a
\ui{Re-check} button --- install a distribution and press it, no restart
needed. Editing works regardless; only typesetting is blocked.
